\section*{Capítulo I, Problema XI}

\textbf{Problema:}

11. Si $A, B$ son dos vectores en el espacio $n$-dimensional, denota por $d(A, B)$ la distancia entre $A$ y $B$, es decir, $d(A, B) = \|B - A\|$. Demuestra que:
\begin{itemize}
    \item $d(A, B) = d(B, A)$,
    \item $d(A, B) \leq d(A, C) + d(B, C)$ para cualesquiera tres vectores $A, B, C$.
\end{itemize}

\textbf{Solución:}

\begin{itemize}
    \item \textbf{Simetría:}

    Por definición, $d(A, B) = \|B - A\|$. Usando la propiedad de la norma, sabemos que $\|B - A\| = \|A - B\|$. Por lo tanto, $d(A, B) = d(B, A)$.

    \item \textbf{Desigualdad triangular:}

    Por la desigualdad triangular de la norma, se cumple que:
    \[
        \|A - C\| + \|C - B\| \geq \|A - B\|.
    \]
    Esto implica que:
    \[
        d(A, C) + d(B, C) \geq d(A, B).
    \]
    Por lo tanto, $d(A, B) \leq d(A, C) + d(B, C)$.
\end{itemize}